% Snapshot of main.tex taken on 2026-11-25: Section 5 built on general theorems (thinning, admissible reserve laws), 16 pages.
% Not part of the build. To compile it, copy it over ../main.tex in the repository root.
\documentclass[format=acmsmall, review=false]{acmart}
\usepackage{acm-ec-26}
\usepackage{booktabs} % For formal tables
\usepackage[ruled]{algorithm2e} % For algorithms
\usepackage{amsthm} % For theorem and remark environments
% amssymb is not loaded: acmart's newtxmath already provides its symbols,
% and loading both fails with "Command \Bbbk already defined".
\usepackage{bbm}
\newcommand{\E}{\mathbb{E}}
\newcommand{\one}{\mathbb{1}}
\renewcommand{\algorithmcfname}{ALGORITHM}
\SetAlFnt{\small}
\SetAlCapFnt{\small}
\SetAlCapNameFnt{\small}
\SetAlCapHSkip{0pt}
\IncMargin{-\parindent}
\newcommand{\OPT}{\mathrm{OPT}}
\newcommand{\SW}{\mathrm{SW}}
\newcommand{\eps}{\varepsilon}

\usepackage{xcolor}

\newcommand{\bruno}[1]{{\color{red}Bruno: #1}}
% Choose a citation style by commenting/uncommenting the appropriate line:
%\setcitestyle{acmnumeric}
\setcitestyle{authoryear}

% Title. Note the optional short title for running heads. In the interest of anonymization, please do not include any acknowledgements.
\title[On sybil-proof combinatorial auctions]{Roses are red, violets are blue, and combinatorial false-name-proof auctions can be efficient too}

% Anonymized submission.
\author{Bruno Mazorra}

% Abstract. Note that this must come before \maketitle.
\begin{abstract}
Iwasaki et al. \cite{iwasaki2010worst} show that, under anonymity, symmetry across goods, and independence of irrelevant goods (IIG), the optimal worst-case welfare approximation ratio of deterministic false-name-proof combinatorial auctions with single-minded bidders is $(m+1)/2$, where $m$ is the number of goods. In this paper, we make two contributions. First, we show that this bound continues to hold without the IIG assumption. Second, we show that allowing randomization substantially improves worst-case welfare. Specifically, we construct a randomized mechanism that is universally truthful, false-name-proof in expectation, runs in polynomial time, and achieves a worst-case welfare approximation ratio of $O(\sqrt{m}\log m)$. Our mechanism incurs only an additional $O(\log m)$ factor in welfare approximation compared to the classical greedy algorithm, showing that randomization makes false-name-proofness compatible with computational efficiency.
\end{abstract}

\begin{document}

% Title page for title and abstract only.
\begin{titlepage}
\maketitle
%\tableofcontents

\end{titlepage}
\section{Introduction}
\section{Related work}
\section{Model}
\label{sec:model}

We consider combinatorial auctions with a single seller, a set
$M=[m]$ of indivisible goods, and a finite set $N$ of bidders,
where $m\ge1$ and $n=|N|$. Bidder $i$ has a private type
$\theta_i=(S_i,v_i)$, consisting of a nonempty target
$S_i\subseteq M$ and a value $v_i\ge0$. With a slight abuse of
notation, we also use $v_i$ for its single-minded valuation
function: $v_i(R)=v_i\mathbf{1}\{S_i\subseteq R\}$ for every
bundle $R\subseteq M$. We denote the type space by $\Theta_M$.
Bidders are risk-neutral and have quasi-linear utilities:
receiving $R$ and paying $P$ gives utility $v_i(R)-P$.


A direct mechanism $\mathcal M=(X,P)$ maps a report profile
$\hat\theta=(\hat\theta_i)_{i\in N}$, with
$\hat\theta_i=(\hat S_i,b_i)\in\Theta_M$, to allocated bundles
$X_i(\hat\theta)\subseteq M$ and payments $P_i(\hat\theta)$
to the seller. The mechanism is defined for every finite set of
participating identities, including the empty set. Allocations
and payments may be randomized; all reports and identities are
fixed before the randomness is realized. Throughout, we require
\emph{ex-post feasibility}: for every report profile,
$X_i(\hat\theta)\cap X_j(\hat\theta)=\varnothing$ almost surely
for distinct $i,j$. Goods may remain unallocated. When no
confusion can arise, we drop the hat on reported targets and
write a report as $(S_i,b_i)$.

Write $u_i(\theta_i;\hat\theta)
=v_i(X_i(\hat\theta))-P_i(\hat\theta)$ for realized utility, and
\[
x_i(\hat\theta)
=\Pr[\hat S_i\subseteq X_i(\hat\theta)],
\qquad
\pi_i(\hat\theta)=\mathbb{E}[P_i(\hat\theta)]
\]
for service probabilities and expected payments. At a truthful
profile $\theta$, bidder $i$'s expected utility is
$U_i(\theta)=\mathbb{E}[u_i(\theta_i;\theta)]
=v_i x_i(\theta)-\pi_i(\theta)$.

We require \emph{individual rationality} (IR) in expectation,
meaning $U_i(\theta)\ge0$ for every bidder and type profile, and
\emph{no positive transfers} (NPT), meaning
$P_i(\hat\theta)\ge0$ almost surely for every report profile.
The expectation in IR is only over the mechanism's randomness,
with all other reports fixed, not over other bidders' types.
Thus $0\le\pi_i(\hat\theta)\le b_i x_i(\hat\theta)$.
For deterministic mechanisms, this is ex-post IR. For randomized
mechanisms, \emph{ex-post IR} additionally requires truthful
realized utility to be nonnegative almost surely.

\subsection{Truthfulness and false-name-proofness}

Let $E$ denote an arbitrary finite profile of external reports,
and write $(\hat\theta_i,E)$ for the profile obtained by adding
identity $i$ with report $\hat\theta_i$. A mechanism is
\emph{truthful} if, for every true type $\theta_i$, alternative
report $\hat\theta_i$, and external profile $E$,
\[
U_i(\theta_i,E)
\ge
\mathbb{E}\bigl[u_i(\theta_i;(\hat\theta_i,E))\bigr].
\]
For deterministic mechanisms, this is dominant-strategy
incentive compatibility (DSIC); for randomized mechanisms, it is
truthfulness in expectation. A deviation may change both the
target bundle and its value. A randomized mechanism is
\emph{universally truthful} if truthful reporting maximizes
realized utility for every realization of its randomness, that
is, if it is a distribution over DSIC mechanisms
\citep{dobzinski2006truthful}.

\emph{False-name-proofness} (FNP) \citep{yokoo2004effect} allows a
physical bidder to
replace its truthful report with any finite family of reports
$\hat\theta_A=(\hat\theta_j)_{j\in A}$ under distinct identities
$A$, its \emph{aliases}, not appearing in
$E$. The bidder values the union of their allocations and pays
the sum of their payments. We require
\[
U_i(\theta_i,E)
\ge
\mathbb{E}\!\left[
  v_i\!\left(\bigcup_{j\in A}X_j(\hat\theta_A,E)\right)
\right]
-\sum_{j\in A}\pi_j(\hat\theta_A,E)
\]
for every true type, external profile, and finite report family.
For randomized mechanisms, this is \emph{FNP in expectation};
for deterministic mechanisms, the expectation disappears.
The family may include the original identity, and different
identities may submit identical reports. Creating identities
is costless, and there is no bound on their number. We allow
$A=\varnothing$, representing nonparticipation, with empty unions
and sums interpreted as $\varnothing$ and zero. All submitted
identities remain in the profile and all assessed payments must
be paid: we consider ordinary FNP, without withdrawal. Taking
$|A|=1$ recovers truthfulness.

\subsection{Anonymity, symmetry, and irrelevant goods}

A mechanism is \emph{anonymous in utilities} if relabelling
bidders leaves their truthful expected utilities unchanged.
For a bijection $\tau:N\to N'$, define
$(\tau\theta)_{\tau(i)}=\theta_i$. Anonymity requires
$U_{\tau(i)}(\tau\theta)=U_i(\theta)$ for every profile, bidder,
and bijection. This definition applies to deterministic and
randomized mechanisms alike, and permits identity-based
tie-breaking when it leaves utilities unchanged.

Separately, a mechanism is \emph{symmetric across goods} if
$U_i(\sigma\theta)=U_i(\theta)$ for every bidder, profile, and
permutation $\sigma$ of $M$. Here $\sigma\theta$ replaces each
type $(S_j,v_j)$ by $(\sigma(S_j),v_j)$, so utilities on the
left are evaluated using the relabelled targets.

\begin{lemma}[Utility symmetry and critical prices]
\label{lem:utility-symmetry}
For a deterministic truthful mechanism satisfying IR and NPT,
anonymity in utilities is equivalent to anonymity of critical
prices. Likewise, symmetry across goods in utilities is
equivalent to symmetry across goods of critical prices.
\end{lemma}

\begin{proof}
Fix an identity $i$, an external profile $E$, and a target $S$.
DSIC makes service monotone in the reported value and gives
all served reports the same payment. IR and NPT imply that
unserved reports pay zero and that the winning payment equals
the critical price $c_{i,E}(S)$, the infimum of winning bids.
Consequently,
\[
U_i((S,v),E)=\max\{0,v-c_{i,E}(S)\}.
\]
If the critical price is infinite, this utility is identically
zero. At a finite critical price, winning and losing both give
zero utility, so tie-breaking does not affect the formula.
The critical price can therefore be recovered from utilities:
\[
c_{i,E}(S)
=\inf\{v\ge0:U_i((S,v),E)>0\},
\]
with $\inf\varnothing=+\infty$. Thus invariance of utilities
under either relabelling implies the corresponding invariance
of critical prices; the converse follows from the utility
formula. Under bidder anonymity, we may consequently write a
common menu $c_E(S)$ depending only on $S$ and the multiset
of external reports. Symmetry across goods further gives
$c_{\sigma E}(\sigma(S))=c_E(S)$.
\end{proof}

For comparison with the deterministic benchmark of
\cite[Definition~2]{iwasaki2010worst}, \emph{independence of
irrelevant goods} (IIG) requires that, for sufficiently large
values, a bidder winning the grand bundle continues to win the
enlarged grand bundle when a new good is added only to its
target, with all values and other targets unchanged.

\subsection{Worst-case welfare}

At a truthful profile $\theta$, realized welfare is
$\mathrm{SW}_{\mathcal M}(\theta)=\sum_{i\in N}v_i(X_i(\theta))$,
with expectation $\sum_{i\in N}v_i x_i(\theta)$, and optimal
welfare is
$\mathrm{OPT}(\theta)=\max_{A\text{ feasible}}\sum_{i\in N}v_i(A_i)$.
For a fixed number $m$ of goods, we define the
\emph{worst-case welfare fraction}\footnote{Its reciprocal is
the usual worst-case welfare approximation ratio, with
$1/0=+\infty$.} by the following infimum over all finite
populations and type profiles on $m$ goods, taking the
expectation over the mechanism's randomness at each fixed
profile:
\[
\rho(\mathcal M)
=
\inf_{\theta:\,\mathrm{OPT}(\theta)>0}
\frac{\mathbb{E}[\mathrm{SW}_{\mathcal M}(\theta)]}
     {\mathrm{OPT}(\theta)}.
\]
\section{Worst-case welfare of deterministic mechanisms}
\label{sec:deterministic}

Throughout this section, the mechanism is deterministic, feasible,
ex-post IR and NPT, DSIC, and ordinary FNP.  It is defined for
every finite population and admits a common bidder-anonymous critical-price
menu, equivariant under permutations of the goods
(Lemma~\ref{lem:utility-symmetry}); as there, we write $c_E(S)$ for this menu at
external profile $E$.  Thus bids strictly above $c_E(S)$ are served, bids
strictly below it are not served, and a winner pays $c_E(S)$; behavior
at equality is fixed by the tie rule.  We do not assume independence of
irrelevant goods.

If $\rho(\mathcal M)>0$, every critical price is finite.  Indeed, if
$c_E(S)=\infty$, then a bidder with target $S$ loses at every finite
value $v$.  The mechanism's welfare is at most $\OPT(E)$, whereas the new
optimum is at least $v$; sending $v\to\infty$ contradicts a positive ratio.
We use this observation without further comment.

The following inequality is the basic deterministic consequence of
false-name-proofness.

\begin{lemma}[False-name cover inequality]
\label{lem:det-cover}
Fix an external profile $E$ and a true target $S$.  Let $\hat\theta_A$ be a
finite false-name family with reports $(S_j,b_j)$, $j\in A$, and let
$D\subseteq A$ be the aliases that win.  If their allocations cover $S$, then
\begin{equation}
  c_E(S)
  \le
  \sum_{j\in D}
  c_{E\cup\hat\theta_{A\setminus\{j\}}}(S_j).
  \label{eq:det-cover}
\end{equation}
Every losing alias remains in the price context on the right-hand side.
\end{lemma}

\begin{proof}
Give the physical bidder a true value $v>c_E(S)$.  Truthful reporting
then gives utility $v-c_E(S)$.  Under the false-name profile, the winning
aliases cover $S$, so the physical value is again $v$.  Losing aliases
pay zero by IR and NPT, while each winner pays the critical price shown
in \eqref{eq:det-cover}.  False-name-proofness lets the common value $v$
cancel and gives the inequality.
\end{proof}

For $h>0$, write $H_h=(M,h)$ for the grand-bundle report of value $h$.
Goods equivariance gives a common clean singleton price
\[
  t(h)=c_{\{H_h\}}(\{g\}),
  \qquad g\in M.
\]

\begin{lemma}[Clean prices have a near-flat pair]
\label{lem:no-iig-flat}
For every $h>0$, $0<t(h)<\infty$.  Moreover, for every $\eps>0$ there exist
$h<h'$ in $[1,2]$ such that
\begin{equation}
  t(h')\le(1+\eps)t(h).
  \label{eq:near-flat}
\end{equation}
\end{lemma}

\begin{proof}
Finiteness was established above.  If $t(h)=0$, then every singleton bid
$\delta>0$ strictly wins against $H_h$.  Feasibility forces the grand-bundle
bidder to lose, so the welfare ratio on this two-bidder profile is at most
$\delta/h$.  Letting $\delta\downarrow0$ contradicts positivity.

For the second claim, fix an integer $\ell\ge1$ and consider the grid
points $1+s/\ell$, $s=0,\ldots,\ell$.  If every adjacent pair satisfied
$t(1+(s+1)/\ell)>(1+\eps)\,t(1+s/\ell)$, then
\[
  t(2)>(1+\eps)^\ell\, t(1).
\]
Because both endpoint prices are positive and finite, this fails for all
sufficiently large $\ell$.  One adjacent pair therefore satisfies
\eqref{eq:near-flat}.
\end{proof}

The next lemma is the replacement for IIG.  A collection of singleton bids
below half the clean price can neither lower another singleton's price nor
co-win with a higher singleton.

\begin{lemma}[Half-price persistence]
\label{lem:no-iig-half}
Fix $h>0$ and a good $g$, write $t=t(h)$, and choose $\lambda<t/2$.  Let $L$
be any family of distinct singleton reports, all of value $\lambda$, on goods
other than $g$.  Then
\begin{equation}
  c_{\{H_h\}\cup L}(\{g\})\ge t.
  \label{eq:no-iig-half}
\end{equation}
Moreover, if $(\{g\},y)$ wins against $\{H_h\}\cup L$ for some $y>\lambda$, no
report in $L$ co-wins.
\end{lemma}

\begin{proof}
We induct on $|L|$, proving the price bound and the no-co-winner statement
together.  Both are immediate when $L=\varnothing$.

Assume the result for every smaller family.  Suppose that $(\{g\},y)$ and a
low report $(\{g'\},\lambda)\in L$ both win.  Put
\[
  L'=L\setminus\{(\{g'\},\lambda)\},
  \qquad
  \gamma=c_{\{H_h,(\{g\},y)\}\cup L'}(\{g'\}).
\]
Because the low report wins, $\gamma\le\lambda<y$.  Raise its bid from
$\lambda$ to $y$.
Its own critical price remains $\gamma$, so it still wins.  Swapping goods
$g$ and $g'$ fixes $H_h$ and every report in $L'$; goods equivariance
therefore gives
\[
  c_{\{H_h,(\{g'\},y)\}\cup L'}(\{g\})=\gamma.
\]
Thus, against the smaller external profile $\{H_h\}\cup L'$, the two reports
$(\{g\},y)$ and $(\{g'\},y)$ both strictly win and pay $\gamma$ each.  Treat
them as two aliases of one bidder with true target $\{g\}$.  The cover
inequality and the induction hypothesis imply
\[
  t
  \le c_{\{H_h\}\cup L'}(\{g\})
  \le 2\gamma
  \le 2\lambda
  <t,
\]
a contradiction.  Hence a high singleton has no low co-winner.

Now choose
\[
  y>\max\{\lambda,c_{\{H_h\}\cup L}(\{g\})\}.
\]
The high singleton wins, and the statement just proved makes every report in
$L$ lose.  Regard the principal and the reports in $L$ as one false-name
family against $H_h$.  Only the principal pays, so
Lemma~\ref{lem:det-cover} gives
\[
  t=c_{\{H_h\}}(\{g\})
  \le c_{\{H_h\}\cup L}(\{g\}).
\]
This proves the price bound and completes the induction.
\end{proof}

\begin{theorem}[Deterministic upper bound without IIG]
\label{thm:no-iig-upper}
Every mechanism in the deterministic class above satisfies
\[
  \rho(\mathcal M)\le\frac{2}{m+1}.
\]
The bound already holds on single-minded star profiles.
\end{theorem}

\begin{proof}
The claim is immediate if $\rho(\mathcal M)=0$, so assume the ratio is
positive and $m\ge2$.  Fix $\eps>0$ and choose the near-flat pair $h<h'$ from
Lemma~\ref{lem:no-iig-flat}.  Write
\[
  t=t(h),
  \qquad
  t'=t(h'),
  \qquad
  \bar t=\max\{t,t'\}.
\]
Then $\bar t\le(1+\eps)t$.  Choose $a>\bar t$ and define the clean reverse
price
\[
  \eta(a)=c_{\{(\{1\},a)\}}(M).
\]
Because $a>t'$, the singleton $(\{1\},a)$ strictly wins against $H_{h'}$.
The grand-bundle report of value $h'$ therefore loses, so
\begin{equation}
  \eta(a)\ge h'>h.
  \label{eq:reverse-slack}
\end{equation}

Singleton auxiliaries cannot lower this grand-bundle price.  Indeed, for any
family $L$ of singleton reports, choose a grand-bundle bid above the finite
threshold $c_{\{(\{1\},a)\}\cup L}(M)$.  The grand report wins and every
singleton in $L$ loses.  Applying Lemma~\ref{lem:det-cover} to this principal
and the losing auxiliaries, against external report $(\{1\},a)$, gives
\begin{equation}
  c_{\{(\{1\},a)\}\cup L}(M)\ge\eta(a).
  \label{eq:dirty-grand-price}
\end{equation}

Fix $0<\lambda<t/2$ and let
\[
  L=\{(\{g\},\lambda):g=2,\ldots,m\}.
\]
Consider the star profile
\begin{equation}
  \theta^\star=\{H_h,(\{1\},a)\}\cup L.
  \label{eq:no-iig-star}
\end{equation}
By \eqref{eq:reverse-slack} and \eqref{eq:dirty-grand-price}, the grand bidder
faces a price strictly above $h$ and loses.

We next show that the high singleton must win.  Suppose instead that it loses
and that a low singleton $(\{g\},\lambda)$ wins.  Let
$L'=L\setminus\{(\{g\},\lambda)\}$.  Half-price persistence gives
\[
  c_{\{H_h\}\cup L'}(\{g\})\ge t.
\]
Against the external profile $\{H_h\}\cup L'$, let one physical bidder submit
the two reports $(\{g\},\lambda)$ and $(\{1\},a)$.  At $\theta^\star$, the
first wins and the second loses.  The false-name family covers $\{g\}$ and
pays at most $\lambda<t$,
contradicting Lemma~\ref{lem:det-cover}.  Thus no low singleton can win when
the high singleton loses.  Since the grand bidder already loses, this would
leave zero welfare on a positive-value profile, contradicting the assumed
positive ratio.  Hence $(\{1\},a)$ wins.

Finally, Lemma~\ref{lem:no-iig-half} implies that none of the low singletons
co-wins with $(\{1\},a)$.  Therefore
\[
  \SW_{\mathcal M}(\theta^\star)=a,
  \qquad
  \OPT(\theta^\star)\ge a+(m-1)\lambda,
\]
and
\[
  \rho(\mathcal M)
  \le\frac{a}{a+(m-1)\lambda}.
\]
Let $a\downarrow\bar t$ and $\lambda\uparrow t/2$.  Since
$\bar t\le(1+\eps)t$,
\[
  \rho(\mathcal M)
  \le
  \frac{2(1+\eps)}{2(1+\eps)+m-1}.
\]
Letting $\eps\downarrow0$ gives $2/(m+1)$.  For $m=1$, the statement reduces
to the trivial bound $\rho\le1$.
\end{proof}

The bound is tight: the Adaptive Reserve Price mechanism of
\cite{iwasaki2010worst} attains $2/(m+1)$.

\section{A randomized mechanism}
\label{sec:randomized}

Randomization removes the linear barrier of
Section~\ref{sec:deterministic}. Our mechanism is assembled from
three ingredients, the first two of which we state in more
generality than we need.

\begin{itemize}
\item \emph{Thinning} (Section~\ref{sec:rand-thinning}). Deleting
goods at random, independently of the reports, and scoring every
report by its expected surviving value turns any mechanism into
another mechanism with the same incentive properties
(Theorem~\ref{thm:thinning}). Thinning is therefore a free tool
for reshaping the order in which a mechanism ranks reports.
\item \emph{Reserve-greedy mechanisms}
(Section~\ref{sec:rand-reserve}). We run the greedy algorithm of
\citet{lehmann2002truth} and impose on its winners a common
reserve, a random fraction of the largest rejected bid.
Theorem~\ref{thm:reserve} characterizes the reserve laws for
which the resulting mechanism is false-name-proof in expectation;
the uniform law is the extremal one.
\item \emph{Scales} (Section~\ref{sec:rand-mechanism}). The final
mechanism mixes thinned reserve-greedy mechanisms over
logarithmically many deletion rates. Only the welfare guarantee
(Section~\ref{sec:rand-welfare}) depends on this mixture.
\end{itemize}

\subsection{Thinning}
\label{sec:rand-thinning}

Following the terminology of \citet{lehmann2002truth}, call a
mechanism \emph{exact} if every bidder receives either its
reported target or nothing; we also require that bidders who
receive nothing pay zero. A bidder is \emph{served} when it
receives its reported target.

\begin{definition}[Thinning]
\label{def:thinning}
Let \(\mathcal M\) be an exact mechanism, and let
\(\widetilde M\subseteq M\) be a random set of \emph{surviving}
goods, drawn independently of the reports and of the randomness
of \(\mathcal M\), with \emph{survival function}
\[
q(R)=\Pr[R\subseteq\widetilde M]>0
\qquad(R\subseteq M).
\]
For a report profile \(\hat\theta\) with reports \((S_i,b_i)\),
let \(q\hat\theta\) denote the \emph{scored profile}, with
reports \((S_i,q(S_i)b_i)\). The \emph{thinning} of
\(\mathcal M\) by \(\widetilde M\) runs \(\mathcal M\) on
\(q\hat\theta\). Bidder \(i\) receives \(S_i\) precisely when it
is served by \(\mathcal M\) and \(S_i\subseteq\widetilde M\); it
then pays \(P_i/q(S_i)\), where \(P_i\) is its payment in
\(\mathcal M\). Every other bidder receives nothing and pays
zero.
\end{definition}

A thinned mechanism ranks a report by the value it is expected
to realize, so that large targets are discounted. The discount is
real: the goods are actually withheld. This is what makes it
compatible with false names.

\begin{theorem}[Thinning preserves incentives]
\label{thm:thinning}
Let \(\mathcal M\) be an exact mechanism that makes no positive
transfers. Its thinning by \(\widetilde M\) is exact and makes
no positive transfers. Moreover:
\begin{enumerate}
\item[(a)] if \(\mathcal M\) is ex-post feasible, respectively
ex-post individually rational, then so is the thinning;
\item[(b)] if \(\mathcal M\) is truthful, respectively
false-name-proof in expectation, then so is the thinning;
\item[(c)] if \(\mathcal M\) is universally truthful and ex-post
individually rational, then the thinning is universally
truthful;
\item[(d)] for every type profile \(\theta\), the expected
welfare of the thinning at \(\theta\) equals
\(\E[\SW_{\mathcal M}(q\theta)]\), the expected welfare of
\(\mathcal M\) at the scored profile.
\end{enumerate}
\end{theorem}

\begin{proof}
The thinning serves a subset of the bidders served by
\(\mathcal M\) and rescales their payments by positive factors.
It is therefore exact, makes no positive transfers, and is
ex-post feasible whenever \(\mathcal M\) is.

Fix a bidder of type \((S,v)\) and an external profile \(E\),
and put \(z=q(S)v\). We compare this bidder, in the thinning and
facing \(E\), with a \emph{fictitious} bidder of type \((S,z)\)
in \(\mathcal M\), facing \(qE\).

\emph{Truthful reports.} If the bidder reports truthfully,
\(\mathcal M\) is run on the fictitious truthful report against
\(qE\). Let \(P\) be the resulting payment in \(\mathcal M\),
which is zero unless the report is served. The realized utility
in the thinning is
\begin{equation}\label{eq:thin-truthful}
\mathbf{1}\{S\subseteq\widetilde M\}
\left(v\,\mathbf{1}\{\text{served}\}-\frac{P}{q(S)}\right)
=\frac{\mathbf{1}\{S\subseteq\widetilde M\}}{q(S)}
\bigl(z\,\mathbf{1}\{\text{served}\}-P\bigr),
\end{equation}
that is, \(\mathbf{1}\{S\subseteq\widetilde M\}/q(S)\) times the
realized truthful utility of the fictitious bidder. This proves
the second half of~(a). Since \(\widetilde M\) is independent of
\(\mathcal M\), taking expectations shows that the two bidders
have the same truthful expected utility, and that the bidder is
served in the thinning with probability \(q(S)\) times the
service probability of the fictitious bidder. Summing
\(v_i\) times this probability over the bidders of a profile
gives~(d).

\emph{False names.} Let the bidder submit aliases \((S_j,b_j)\),
\(j\in A\), so that \(\mathcal M\) is run on the scored aliases
\((S_j,q(S_j)b_j)\) against \(qE\), and let \(D\subseteq A\) be
the random set of aliases served by \(\mathcal M\). No deleted
good is ever allocated. The bidder therefore obtains \(S\) in the
thinning only if \(S\subseteq\bigcup_{j\in D}S_j\) and
\(S\subseteq\widetilde M\), and by independence its expected
value is at most
\[
v\,q(S)\Pr\Bigl[S\subseteq\bigcup_{j\in D}S_j\Bigr]
=z\Pr\Bigl[S\subseteq\bigcup_{j\in D}S_j\Bigr],
\]
which is the expected value that the fictitious bidder derives
from the scored aliases in \(\mathcal M\). Alias \(j\) pays
\(P_j/q(S_j)\) when \(j\in D\) and \(S_j\subseteq\widetilde M\),
so its expected payment in the thinning is
\(\E[P_j\mathbf{1}\{j\in D\}]=\pi_j\), its expected payment in
\(\mathcal M\). The deviation's expected utility in the thinning
is thus at most the expected utility that the fictitious bidder
derives from the scored aliases. If \(\mathcal M\) is
false-name-proof in expectation, this is at most the fictitious
truthful expected utility, which equals the truthful expected
utility in the thinning. The case \(|A|=1\) is truthfulness.
This proves~(b).

\emph{Universal truthfulness.} Fix a realization of the
randomness of \(\mathcal M\) and of \(\widetilde M\), and let the
bidder report \((R,b)\). If this report is not served in the
thinning, or if \(S\not\subseteq R\), the realized utility is at
most zero, hence at most the truthful one by~(a). Otherwise
\(S\subseteq R\subseteq\widetilde M\), the report
\((R,q(R)b)\) is served by \(\mathcal M\) at some payment
\(P\ge0\), and since \(q(R)\le q(S)\) the realized utility is
\[
v-\frac{P}{q(R)}\le v-\frac{P}{q(S)}=\frac{z-P}{q(S)}.
\]
Here \(z-P\) is the realized utility in \(\mathcal M\) of the
fictitious bidder when it reports \((R,q(R)b)\). By universal
truthfulness of \(\mathcal M\), it is at most the fictitious
truthful realized utility, and~\eqref{eq:thin-truthful}
gives~(c), because \(S\) survives.
\end{proof}

\subsection{Reserve-greedy mechanisms}
\label{sec:rand-reserve}

\paragraph{Greedy notation.}
A \emph{weighted report} \((S_i,w_i)\) consists of a target and a
weight \(w_i>0\). For an indexed collection \(T\) of weighted
reports, let \(G(T)\) be the packing obtained by processing
reports in nonincreasing order of weight and accepting a report
whenever its target is disjoint from the previously accepted
targets. Ties are resolved by a priority order \(\omega\) on the
identities, independent of the reported values and targets, and
every subprofile inherits the induced order. Define
\[
C(T)=\max_{i\in T\setminus G(T)}w_i,
\]
the largest rejected weight, where the maximum of an empty set is
zero. For a weighted profile \(E\) and a nonempty bundle
\(R\subseteq M\), define
\[
\kappa(R;E)=
\max\bigl\{
w_j:j\in G(E),\
S_j\cap R\ne\varnothing
\bigr\},
\qquad
B(R;E)=\max\{C(E),\kappa(R;E)\},
\]
again using zero for an empty maximum. The \emph{greedy critical
weight} of \(i\in G(T)\) is
\(\kappa_i(T)=\kappa(S_i;T\setminus\{i\})\).

\begin{definition}[Reserve-greedy mechanism \(\mathcal A^{\mu}\)]
\label{def:reserve-greedy}
Let \(\mu\) be a probability law on \([0,1]\), the
\emph{reserve law}. Reports with value zero are discarded. Given
the remaining reports \((S_i,b_i)\), proceed as follows.
\begin{enumerate}
\item \emph{Greedy.} Let \(T\) be the weighted profile with
weights \(w_i=b_i\). Draw a uniformly random priority order
\(\omega\) on the participating identities, independently of the
reports, and compute
\[
G=G(T),\qquad
C=C(T),\qquad
\kappa_i=\kappa_i(T)\quad(i\in G).
\]
\item \emph{Reserve.} Draw \(\xi\sim\mu\).
\item \emph{Allocation and payments.} Bidder \(i\) receives
\(S_i\) precisely when \(i\in G\) and \(w_i\ge\xi C\), and then
pays
\begin{equation}\label{eq:base-payment}
P_i=\max\{\kappa_i,\xi C\}.
\end{equation}
Every other bidder receives nothing and pays zero.
\end{enumerate}
\end{definition}

If \(\xi\) is identically zero, \(\mathcal A^{\mu}\) is the
greedy mechanism with critical payments of
\citet{lehmann2002truth}, ranked by value. That mechanism is
truthful but not false-name-proof: two aliases can shield each
other from a common competitor, which reduces both critical
weights to zero. The reserve \(\xi C\) charges for this
shielding. It is computed from rejected reports, which a bidder
can create under false names but can never remove.

\begin{theorem}[Admissible reserve laws]
\label{thm:reserve}
For every reserve law \(\mu\), the mechanism \(\mathcal A^{\mu}\)
is exact, ex-post feasible, universally truthful, ex-post
individually rational, and makes no positive transfers.
Furthermore, let \(\xi\sim\mu\) and consider the condition
\begin{equation}\label{eq:R}
\tag{R}
\E\bigl[\xi\,\mathbf{1}\{\xi\le\beta\}\bigr]
\ge\frac{\beta}{2}\Pr[\xi\le\beta]
\qquad\text{for every }\beta\in[0,1].
\end{equation}
\begin{enumerate}
\item[(a)] If \(\mu\) satisfies~\eqref{eq:R}, then
\(\mathcal A^{\mu}\) is false-name-proof in expectation.
\item[(b)] Conversely, if \(m\ge3\) and \(\mathcal A^{\mu}\) is
false-name-proof in expectation, then \(\mu\)
satisfies~\eqref{eq:R}.
\end{enumerate}
\end{theorem}

Condition~\eqref{eq:R} requires that, conditional on lying below
any level \(\beta\), the reserve fraction is on average at least
\(\beta/2\): two aliases that shield each other at level
\(\beta\) must together pay what a single bidder would. The proof
is given in Section~\ref{sec:rand-reserve-proof}. Three examples
illustrate the condition.

\begin{itemize}
\item The uniform law on \([0,1]\) satisfies~\eqref{eq:R} with
equality at every level, and it is the only law that does.
Indeed, after integration by parts, equality at every \(\beta\)
reads
\(\beta\Pr[\xi\le\beta]=2\int_0^\beta\Pr[\xi\le\zeta]\,d\zeta\),
whose only solution with \(\Pr[\xi\le1]=1\) is
\(\Pr[\xi\le\beta]=\beta\). More generally, \eqref{eq:R} is
equivalent to
\(\beta\Pr[\xi\le\beta]\ge2\int_0^\beta\Pr[\xi\le\zeta]\,d\zeta\),
which makes \(\beta^{-2}\int_0^\beta\Pr[\xi\le\zeta]\,d\zeta\)
nondecreasing in \(\beta\). Every admissible law therefore
satisfies \(\int_0^\beta\Pr[\xi\le\zeta]\,d\zeta\le\beta^2/2\),
with equality for the uniform law: averaged over levels, no
admissible reserve serves low bids more often.
\item The law concentrated at zero violates~\eqref{eq:R}, in
accordance with the shielding attack described above.
\item A point mass at \(\beta_0\ge1/2\) satisfies~\eqref{eq:R}.
With a fixed priority order it gives a deterministic
false-name-proof mechanism, in which greedy winners below the
reserve \(\beta_0C\) are excluded. Its welfare is poor, as
Section~\ref{sec:deterministic} leads one to expect: two unit
bids on one good exclude bids just below \(1/2\) on each of the
other \(m-1\) goods, so its approximation ratio is at least
\((m+1)/2\).
\end{itemize}
Randomizing the reserve is thus what allows winners with low
bids to be served at all, and~\eqref{eq:R} determines how often.

\subsection{The mechanism and the main result}
\label{sec:rand-mechanism}

Let
\[
\mathcal K=\bigl\{1,2,4,\ldots,2^{\lfloor\log_2m\rfloor}\bigr\},
\qquad
|\mathcal K|=1+\lfloor\log_2m\rfloor.
\]
A \emph{scale} is an element \(k\in\mathcal K\cup\{\infty\}\). At
scale \(k\), every good survives independently with probability
\[
p_k=1-\frac{1}{2k},
\qquad
p_\infty=1,
\]
so that a bundle \(R\) survives with probability
\(q_k(R)=p_k^{|R|}\).

\begin{definition}[The mechanisms \(\mathcal A_k\) and
\(\mathcal A_{\mathcal K}\)]
\label{def:M}
Let \(\mathcal A_\infty\) be the reserve-greedy mechanism whose
reserve law is uniform on \([0,1]\). For \(k\in\mathcal K\), let
\(\mathcal A_k\) be the thinning of \(\mathcal A_\infty\) by a
random set \(\widetilde M\) that retains each good independently
with probability \(p_k\). The mechanism
\(\mathcal A_{\mathcal K}\) runs \(\mathcal A_\infty\) with
probability \(1/2\); with probability \(1/2\), it chooses a
scale \(k\) uniformly from \(\mathcal K\) and runs
\(\mathcal A_k\).
\end{definition}

Unfolding the definitions, \(\mathcal A_k\) gives report \(i\)
the weight \(w_i=q_k(S_i)b_i\), runs greedy on the resulting
weighted profile \(T\), and draws
\(\xi\sim\operatorname{Unif}[0,1]\) and \(\widetilde M\). Bidder
\(i\) receives \(S_i\) precisely when
\[
i\in G,\qquad w_i\ge\xi C,\qquad S_i\subseteq\widetilde M,
\]
and then pays
\begin{equation}\label{eq:actual-payment}
P_i=\frac{\max\{\kappa_i,\xi C\}}{q_k(S_i)}.
\end{equation}

\begin{theorem}\label{thm:main}
For every \(m\ge1\), the mechanism \(\mathcal A_{\mathcal K}\)
for single-minded bidders is universally truthful,
false-name-proof in expectation, ex-post individually rational,
and makes no positive transfers. For every type profile
\(\theta\),
\[
\E[\SW_{\mathcal A_{\mathcal K}}(\theta)]
\ge
\frac{\OPT(\theta)}{
4\sqrt{2m\ln(em)\bigl(1+\lfloor\log_2m\rfloor\bigr)}+8/e}.
\]
In particular, its worst-case welfare approximation ratio
\(1/\rho(\mathcal A_{\mathcal K})\) is
\[
O(\sqrt m\log(m+1))=\widetilde O(\sqrt m).
\]

The mechanism is anonymous in utilities and symmetric across
goods, and it uses polynomially many arithmetic operations. No
independence-of-irrelevant-goods condition is imposed.
\end{theorem}

The incentive properties follow at once from the two general
theorems.

\begin{corollary}\label{cor:mixture}
Every mechanism that draws a scale
\(k\in\mathcal K\cup\{\infty\}\) from a fixed distribution,
independently of the reports, and then runs \(\mathcal A_k\) is
ex-post feasible, universally truthful, ex-post individually
rational, false-name-proof in expectation, and makes no positive
transfers. In particular, this holds for
\(\mathcal A_{\mathcal K}\).
\end{corollary}

\begin{proof}
The uniform law satisfies~\eqref{eq:R}, so \(\mathcal A_\infty\)
has all five properties by Theorem~\ref{thm:reserve}. Each
\(\mathcal A_k\) inherits them by Theorem~\ref{thm:thinning}.
The ex-post properties then hold for every realization of the
scale. Because the scale is drawn independently of the reports,
the expected utility of any family of reports is the average
over the scale of its expected utility in \(\mathcal A_k\), and
the false-name inequality is preserved by this average.
\end{proof}

The welfare bound is Proposition~\ref{prop:welfare}, and the
remaining claims of Theorem~\ref{thm:main} are
Proposition~\ref{prop:implementation}.

\subsection{Proof of Theorem~\ref{thm:reserve}}
\label{sec:rand-reserve-proof}

\begin{lemma}[Greedy normalizer]
\label{lem:greedy}
Let \(T\) and \(E\) be weighted profiles.
\begin{enumerate}
\item[(a)] We have
\begin{equation}\label{eq:pair}
C(T)=
\max_{\substack{i,j\in T,\ i\ne j\\
                 S_i\cap S_j\ne\varnothing}}
\min\{w_i,w_j\}.
\end{equation}
In particular, \(C(T)\) does not depend on the priority order and
cannot decrease when reports are added.
\item[(b)] Add to \(E\) a report \((R,w)\) under a new identity.
It is accepted by greedy if \(w>\kappa(R;E)\) and rejected if
\(w<\kappa(R;E)\). Whenever it is accepted,
\begin{equation}\label{eq:normalizer}
C\bigl(E\cup\{(R,w)\}\bigr)=B(R;E).
\end{equation}
\end{enumerate}
\end{lemma}

\begin{proof}
(a) At least one member of every conflicting pair is rejected,
and the rejected member has weight at most \(C(T)\); hence the
minimum weight of the pair is at most \(C(T)\). Conversely, every
rejected report is blocked by an accepted report of at least its
weight, and the minimum weight of this pair is the rejected
weight. This gives the reverse inequality.

(b) Reports processed before \((R,w)\) are treated exactly as in
the greedy run on \(E\). The new report is therefore accepted
precisely when no report of \(G(E)\) processed before it has a
target intersecting \(R\). If \(w>\kappa(R;E)\), every such
report has smaller weight and is processed later. If
\(w<\kappa(R;E)\), a report of \(G(E)\) of weight
\(\kappa(R;E)\) whose target intersects \(R\) is processed
earlier.

Suppose now that the new report is accepted, so that
\(w\ge\kappa(R;E)\). A report of \(E\) that conflicts with it is
either accepted in \(G(E)\), and therefore has weight at most
\(\kappa(R;E)\), or rejected there, and therefore has weight at
most \(C(E)\). Conflicting pairs inside \(E\) have minimum weight
at most \(C(E)\). Equation~\eqref{eq:pair} gives the upper bound
\(C(E\cup\{(R,w)\})\le B(R;E)\). The lower bound follows from
monotonicity and, when \(\kappa(R;E)>0\), from the conflicting
report of \(G(E)\) of weight \(\kappa(R;E)\le w\).
\end{proof}

\begin{lemma}[Realized critical prices]
\label{lem:threshold}
Consider \(\mathcal A^{\mu}\). Fix a realization of
\((\omega,\xi)\) and an external profile \(E\). For a nonempty
bundle \(R\subseteq M\), let
\[
c_E(R)=\max\{\kappa(R;E),\,\xi B(R;E)\}.
\]
A bidder reporting \((R,b)\) against \(E\) is served and pays
\(c_E(R)\) if \(b>c_E(R)\), and is not served if \(b<c_E(R)\).
If \(b=c_E(R)\), it is either served at payment \(b\) or not
served. Moreover, \(c_E(R)\le c_E(R')\) whenever
\(R\subseteq R'\).
\end{lemma}

\begin{proof}
By Lemma~\ref{lem:greedy}(b), the report is accepted if
\(b>\kappa(R;E)\) and rejected if \(b<\kappa(R;E)\); whenever it
is accepted, its greedy critical weight is \(\kappa(R;E)\) and
the normalizer of the full profile is \(B(R;E)\). An accepted
report is thus served precisely when \(b\ge\xi B(R;E)\), and it
then pays \(c_E(R)\) by~\eqref{eq:base-payment}. Consequently,
the report is served if \(b>c_E(R)\), it is not served if
\(b<c_E(R)\), and at equality it pays \(c_E(R)=b\) if it is
served. Monotonicity holds because \(\kappa(R;E)\), and hence
\(B(R;E)\), can only increase when \(R\) is enlarged.
\end{proof}

\begin{proof}[Proof of Theorem~\ref{thm:reserve}, ex-post
properties]
The mechanism is exact by construction, its served bidders form
a subset of a greedy packing, and payments are nonnegative. By
Lemma~\ref{lem:threshold}, every realization of
\(\mathcal A^{\mu}\) is a deterministic mechanism in which a
report wins when its value exceeds a critical price that does
not depend on that value and is nondecreasing in the reported
target, and in which winners pay the critical price. Such
mechanisms are truthful and individually rational for
single-minded bidders
\citep{lehmann2002truth,mualem2008truthful}; in the terminology
of \citet{yokoo2003characterization}, they are price-oriented,
rationing-free protocols. For completeness: the truthful
realized utility of a bidder of type \((S,v)\) is
\(\max\{0,v-c_E(S)\}\ge0\); a report \((R,b)\) that is served
with \(S\subseteq R\) yields \(v-c_E(R)\le v-c_E(S)\); and every
other report yields at most zero.
\end{proof}

The next lemma is the combinatorial core of the false-name
argument. It concerns greedy packing only, and expresses that a
competitor displaced by false names leaves a trace.

\begin{lemma}[Displacement]
\label{lem:displacement}
Let \(E\) be a weighted profile, let \(S\subseteq M\) be a
nonempty bundle, and put \(\kappa=\kappa(S;E)\) and
\(B=B(S;E)\). Let \(T'\) be obtained from \(E\) by adding
finitely many weighted reports under new identities \(A\), with
a priority order extending that of \(E\). Let \(F\subseteq A\)
be the set of new reports accepted by \(G(T')\), and suppose
that the targets of \(F\) cover \(S\). Then
\begin{equation}\label{eq:C-comparison}
C(T')\ge B.
\end{equation}
Moreover, if \(\kappa>0\), then one of the following holds:
\begin{enumerate}
\item[(a)] at least two reports of \(F\) have weight at least
\(\kappa\);
\item[(b)] exactly one report \(j\in F\) has weight at least
\(\kappa\), and \(\kappa_j(T')\ge\kappa\).
\end{enumerate}
\end{lemma}

\begin{proof}
By Lemma~\ref{lem:greedy}(a), \(C(T')\ge C(E)\). If \(\kappa>0\),
fix a report of \(G(E)\) of weight \(\kappa\) whose target
intersects \(S\), and call it the \emph{blocker}. The blocker
cannot be accepted in \(G(T')\), because the targets of \(F\)
cover \(S\) and a greedy packing is disjoint. It is therefore
rejected in \(T'\), implying \(C(T')\ge\kappa\). This
proves~\eqref{eq:C-comparison}.

Suppose now that \(\kappa>0\) and that~(a) fails. Since \(F\) is
nonempty, there is \(j\in F\) such that every report of \(F\)
other than \(j\) has weight below \(\kappa\). We show that
\(\kappa_j(T')\ge\kappa\). Because \(j\) is accepted,
Lemma~\ref{lem:greedy}(b) then gives
\(w_j\ge\kappa_j(T')\ge\kappa\), which is~(b).

Suppose, to the contrary, that \(\kappa_j(T')<\kappa\): in the
profile \(T'\setminus\{j\}\), no accepted report of weight at
least \(\kappa\) has a target intersecting \(S_j\). Call the
reports of weight at least \(\kappa\) the \emph{prefix}; greedy
processes them first. We claim that \(G(T')\) and
\(G(T'\setminus\{j\})\) accept the same prefix reports other
than \(j\). This is automatic if \(w_j<\kappa\). Otherwise, scan
the prefix in greedy order. The two runs coincide before \(j\),
and \(j\) is accepted in \(G(T')\). A later prefix report
accepted in \(G(T'\setminus\{j\})\) has a target disjoint from
\(S_j\) by assumption, so the presence of \(j\) never changes a
decision.

By the claim, a new report in the prefix that is accepted in
\(G(T'\setminus\{j\})\) belongs to \(F\setminus\{j\}\) and has
weight at least \(\kappa\); no such report exists. Hence all new
reports in the prefix are rejected in \(G(T'\setminus\{j\})\).
Deleting rejected reports never changes a greedy outcome, so the
prefix reports accepted in \(G(T'\setminus\{j\})\) are exactly
the prefix reports of \(G(E)\). In particular, the blocker is
accepted in \(G(T'\setminus\{j\})\) and, by the claim, in
\(G(T')\). This contradicts the first paragraph of the proof.
\end{proof}

\begin{proof}[Proof of Theorem~\ref{thm:reserve}(a)]
Fix a bidder of type \((S,v)\), the external profile \(E\), and
the priority order among the external identities. Put
\(\kappa=\kappa(S;E)\) and \(B=B(S;E)\). By
Lemma~\ref{lem:threshold}, the truthful realized utility is
\(\max\{0,v-\max\{\kappa,\xi B\}\}\), whatever the bidder's own
priority.

Replace the truthful report by any finite family
\(\hat\theta_A\) of alias reports \((S_j,w_j)\), with arbitrary
targets, values, and priorities. Let \(T'\) be the resulting
profile, let \(F\subseteq A\) be the set of aliases accepted by
\(G(T')\), and write \(C'=C(T')\). Abstention gives zero utility.
If the targets of \(F\) do not cover \(S\), the deviation
produces no value and has nonnegative payments. In both cases it
cannot improve upon truthful utility. Suppose henceforth that
the targets of \(F\) cover \(S\), and define
\[
F_S=\{j\in F:S_j\cap S\ne\varnothing\},
\qquad
\alpha=\min_{j\in F_S}\min\{w_j,C'\},
\]
with the minimum attained by \(j_0\in F_S\).

Fix \(\xi\). Since \(\xi\le1\), an alias \(j\in F\) is served
precisely when \(\xi C'\le\min\{w_j,C'\}\), and it then pays
\(\max\{\kappa_j(T'),\xi C'\}\); the other aliases pay nothing.
Every alias in \(F_S\) is necessary to supply \(S\), because the
accepted targets are disjoint, so the bidder obtains \(S\) only
if \(\xi C'\le\alpha\). Its realized utility is therefore at
most
\begin{equation}\label{eq:attack}
v\,\mathbf{1}\{\xi C'\le\alpha\}
-\sum_{j\in F}
\mathbf{1}\{\xi C'\le\min\{w_j,C'\}\}
\max\{\kappa_j(T'),\xi C'\}.
\end{equation}
We claim the following \emph{cover inequality in expectation}:
\begin{equation}\label{eq:cover}
\E_\xi\Bigl[\,
\sum_{j\in F}
\mathbf{1}\{\xi C'\le\min\{w_j,C'\}\}
\max\{\kappa_j(T'),\xi C'\}
\Bigr]
\ge
\E_\xi\bigl[
\mathbf{1}\{\xi C'\le\alpha\}\max\{\kappa,\xi C'\}
\bigr].
\end{equation}
It is the randomized counterpart of
Lemma~\ref{lem:det-cover}: the aliases together pay, on average,
at least what a single report served on the same event would pay
against the competitor of weight \(\kappa\).
Granting~\eqref{eq:cover}, and using \(C'\ge B\)
from~\eqref{eq:C-comparison},
\[
\E[u_{\mathrm{aliases}}]
\le
\E_\xi\bigl[
\mathbf{1}\{\xi C'\le\alpha\}
\bigl(v-\max\{\kappa,\xi C'\}\bigr)
\bigr]
\le
\E_\xi\bigl[\max\{0,v-\max\{\kappa,\xi B\}\}\bigr]
=
\E[u_{\mathrm{truthful}}].
\]
The argument is conditional on the external priorities and
permits arbitrary priorities for the aliases. It therefore
remains valid after averaging over the uniformly random priority
order.

It remains to prove~\eqref{eq:cover}. We compare the payments of
suitable aliases with the right-hand integrand, in three cases.

\emph{Case 1: \(\kappa=0\).} For every \(\xi\), the payment of
\(j_0\) is at least
\(\mathbf{1}\{\xi C'\le\alpha\}\,\xi C'\), which is the
right-hand integrand.

\emph{Case 2: \(\kappa>0\) and
Lemma~\ref{lem:displacement}(b) holds.} Let \(j\) be the unique
alias of \(F\) of weight at least \(\kappa\), so that
\(\kappa_j(T')\ge\kappa\). Then \(\min\{w_j,C'\}\ge\alpha\): this
is an equality if \(j=j_0\), and otherwise
\(\alpha\le w_{j_0}<\kappa\le\min\{w_j,C'\}\). Hence, for every
\(\xi\), the payment of \(j\) is at least
\(\mathbf{1}\{\xi C'\le\alpha\}\max\{\kappa,\xi C'\}\).

\emph{Case 3: \(\kappa>0\) and
Lemma~\ref{lem:displacement}(a) holds.} Call the aliases of
\(F\) of weight at least \(\kappa\) \emph{heavy}; there are at
least two. Put \(\beta=\kappa/C'\in(0,1]\). A heavy alias is
served whenever \(\xi\le\beta\), and then pays at least
\(\xi C'\). If \(\alpha\ge\kappa\), let \(j\ne j_0\) be heavy.
The payments of \(j_0\) and \(j\) total at least
\[
\mathbf{1}\{\xi C'\le\alpha\}\,\xi C'
+\mathbf{1}\{\xi\le\beta\}\,\xi C',
\]
whereas
\[
\mathbf{1}\{\xi C'\le\alpha\}\max\{\kappa,\xi C'\}
=\mathbf{1}\{\xi\le\beta\}\,\kappa
+\mathbf{1}\{\kappa<\xi C'\le\alpha\}\,\xi C'.
\]
If \(\alpha<\kappa\), then \(j_0\) is not heavy, two heavy
aliases pay at least
\(2\cdot\mathbf{1}\{\xi\le\beta\}\,\xi C'\) in total, and
\(\mathbf{1}\{\xi C'\le\alpha\}\max\{\kappa,\xi C'\}
\le\mathbf{1}\{\xi\le\beta\}\,\kappa\). In both subcases, the
total payment exceeds the right-hand integrand by at least
\[
\mathbf{1}\{\xi\le\beta\}\,(2\xi-\beta)\,C',
\]
whose expectation is nonnegative by~\eqref{eq:R}.
\end{proof}

\begin{proof}[Proof of Theorem~\ref{thm:reserve}(b)]
Suppose that~\eqref{eq:R} fails at some \(\beta\in(0,1]\).
Consider the goods \(1,2,3\) and the external profile \(E\)
consisting of two reports with target \(\{3\}\) and value \(1\),
and one report with target \(\{1,2\}\) and value \(\beta\). Then
\(C(E)=1\), \(\kappa(\{1,2\};E)=\beta\), and
\(B(\{1,2\};E)=1\), so a bidder of type \((\{1,2\},\beta)\) has
truthful realized utility
\(\max\{0,\beta-\max\{\beta,\xi\}\}=0\).

Let this bidder instead submit two aliases with targets
\(\{1\}\) and \(\{2\}\) and common value \(\beta'>\beta\). Greedy
accepts both aliases and rejects the report with target
\(\{1,2\}\). The normalizer is still \(1\), and both aliases
have greedy critical weight zero, each shielding the other. Both
are therefore served when \(\xi\le\beta'\), and each then pays
\(\xi\). The deviation's expected utility is
\[
\beta\Pr[\xi\le\beta']
-2\,\E\bigl[\xi\,\mathbf{1}\{\xi\le\beta'\}\bigr].
\]
As \(\beta'\) decreases to \(\beta\), this converges to
\(\beta\Pr[\xi\le\beta]-2\,\E[\xi\,\mathbf{1}\{\xi\le\beta\}]>0\).
The deviation is thus profitable for \(\beta'\) close to
\(\beta\), and \(\mathcal A^{\mu}\) is not false-name-proof in
expectation.
\end{proof}

\subsection{Welfare of the mixture}
\label{sec:rand-welfare}

The welfare analysis combines the charging argument of
\citet{lehmann2002truth} for greedy with a second-moment bound
that accounts for the reserve.

\begin{lemma}[Welfare of \(\mathcal A_\infty\)]
\label{lem:welfare-base}
Let \(\theta\) be a type profile and let \(V\ge\max_iv_i\). Run
\(\mathcal A_\infty\) on \(\theta\) with a fixed priority order
\(\omega\), let \(G\) and \(C\) be the resulting greedy packing
and normalizer, and let \(\varphi\) be the expected welfare over
\(\xi\).
\begin{enumerate}
\item[(a)] We have
\(\varphi=\sum_{i\in G}v_i\min\{1,v_i/C\}\), with the minimum
replaced by one if \(C=0\).
\item[(b)] We have \(\varphi\ge\max_iv_i\) and
\(\sum_{i\in G}v_i^2\le V\varphi\).
\item[(c)] For every set \(I\subseteq N\) of bidders with
pairwise disjoint targets and every \(\chi\ge1\),
\[
\sum_{i\in I}v_i
\le
\sqrt{\chi\,|I|\,V\varphi}
+|I|\max\{v_j:j\in G,\ |S_j|>\chi\},
\]
where the maximum of an empty set is zero.
\end{enumerate}
\end{lemma}

\begin{proof}
(a) Weights equal values, and \(i\in G\) is served when
\(v_i\ge\xi C\), which has probability \(\min\{1,v_i/C\}\).

(b) The report processed first by greedy has the largest value,
is accepted, and has weight at least \(C\); it is served for
every \(\xi\). Moreover, \(v_i\le V\) and \(C\le V\) give
\(v_i^2\le Vv_i\min\{1,v_i/C\}\).

(c) Assign each bidder of \(I\) to the first greedy winner that
blocks it, or to itself if it is accepted, and let \(\nu_j\) be
the number of bidders assigned to \(j\in G\). Greedy ordering
and disjointness of the targets in \(I\) imply
\[
\sum_{i\in I}v_i\le\sum_{j\in G}\nu_jv_j,
\qquad
\nu_j\le|S_j|,
\qquad
\sum_{j\in G}\nu_j=|I|.
\]
For winners with \(|S_j|\le\chi\), we have
\(\sum\nu_j^2\le\chi\,|I|\), and Cauchy--Schwarz with~(b) gives
\[
\sum_{\substack{j\in G\\|S_j|\le\chi}}\nu_jv_j
\le\sqrt{\chi\,|I|\,V\varphi}.
\]
The remaining winners contribute at most \(|I|\) times their
largest value.
\end{proof}

\begin{proposition}[Welfare guarantee]
\label{prop:welfare}
For every type profile \(\theta\),
\[
\E[\SW_{\mathcal A_{\mathcal K}}(\theta)]
\ge
\frac{\OPT(\theta)}{
4\sqrt{2m\ln(em)\,|\mathcal K|}+8/e}.
\]
\end{proposition}

\begin{proof}
Put \(V=\max_iv_i\). If \(V=0\), the claim is immediate.

For a scale \(k\in\mathcal K\cup\{\infty\}\), let \(\Phi_k\)
denote the expected welfare of \(\mathcal A_k\) at \(\theta\). By
Theorem~\ref{thm:thinning}(d), \(\Phi_k\) is the expected welfare
of \(\mathcal A_\infty\) at the scored profile \(q_k\theta\),
whose values \(w_i=q_k(S_i)v_i\) are at most \(V\). For a fixed
priority order \(\omega\), let \(\varphi_{k,\omega}\) denote the
expected welfare of \(\mathcal A_\infty\) at \(q_k\theta\) over
\(\xi\), so that \(\Phi_k=\E_\omega[\varphi_{k,\omega}]\).

\paragraph{The scale \(\infty\).}
By Lemma~\ref{lem:welfare-base}(b),
\begin{equation}\label{eq:top}
\Phi_\infty\ge V.
\end{equation}

\paragraph{The scales in \(\mathcal K\).}
Fix an optimal packing \(N^\star\subseteq N\), that is, a set of
bidders with pairwise disjoint targets, positive values, and
total value \(\OPT(\theta)\). For each \(k\in\mathcal K\), let
\[
N_k=\{i\in N^\star:k\le |S_i|<2k\},
\]
and define
\[
W_k=\sum_{i\in N_k}v_i,\qquad
n_k=|N_k|,\qquad
m_k=\sum_{i\in N_k}|S_i|.
\]
Every target has between \(1\) and \(m\) goods, so the classes
\(N_k\) partition \(N^\star\). Disjointness gives
\begin{equation}\label{eq:mass}
k\,n_k\le m_k,\qquad
\sum_{k\in\mathcal K}m_k\le m,\qquad
\sum_{k\in\mathcal K}n_k\le m.
\end{equation}
For \(i\in N_k\), we have \(|S_i|<2k\), and
\((1-1/(2k))^{2k}\) is nondecreasing in \(k\ge1\), so
\[
q_k(S_i)
\ge
\left(1-\frac{1}{2k}\right)^{2k}
\ge\left(1-\frac12\right)^{2}
=\frac14,
\]
and hence \(v_i\le4w_i\).

Set \(d=\ln(em)\), and apply Lemma~\ref{lem:welfare-base}(c) to
the profile \(q_k\theta\), with \(I=N_k\) and \(\chi=2kd\). A
greedy winner \(j\) with \(|S_j|>2kd\) has
\[
w_j\le q_k(S_j)V
\le e^{-|S_j|/(2k)}V
<\frac{V}{em}.
\]
Therefore, using~\eqref{eq:mass},
\[
\frac{W_k}{4}
\le\sum_{i\in N_k}w_i
\le\sqrt{2kd\,n_kV\varphi_{k,\omega}}+\frac{Vn_k}{em}
\le\sqrt{2d\,m_kV\varphi_{k,\omega}}+\frac{Vn_k}{em}.
\]
Averaging over \(\omega\) and using concavity of the square root
gives
\begin{equation}\label{eq:class}
W_k
\le
4\sqrt{2d\,m_kV\Phi_k}
+
\frac{4Vn_k}{em}.
\end{equation}

Sum~\eqref{eq:class} over \(k\in\mathcal K\) and apply
Cauchy--Schwarz and~\eqref{eq:mass}:
\begin{equation}\label{eq:opt}
\begin{aligned}
\OPT(\theta)
&\le
4\sqrt{2dV}\sum_{k\in\mathcal K}\sqrt{m_k\Phi_k}
+\frac{4V}{e}\\
&\le
4\sqrt{2mdV\sum_{k\in\mathcal K}\Phi_k}
+\frac{4V}{e}.
\end{aligned}
\end{equation}

\paragraph{The mixture.}
By Definition~\ref{def:M} and~\eqref{eq:top}, the expected
welfare \(\Phi\) of \(\mathcal A_{\mathcal K}\) satisfies
\[
\Phi
=\frac12\Phi_\infty
+\frac1{2|\mathcal K|}\sum_{k\in\mathcal K}\Phi_k
\ge\frac12V
+\frac1{2|\mathcal K|}\sum_{k\in\mathcal K}\Phi_k.
\]
The arithmetic--geometric mean inequality yields
\[
V\sum_{k\in\mathcal K}\Phi_k\le|\mathcal K|\,\Phi^2,
\qquad
V\le2\Phi.
\]
Substituting into~\eqref{eq:opt} proves
\[
\OPT(\theta)
\le
\left(
4\sqrt{2m\ln(em)\,|\mathcal K|}+\frac8e
\right)\Phi,
\]
which is the claimed approximation guarantee.
\end{proof}

The proof isolates the role of each ingredient. The reserve
costs a second moment, Lemma~\ref{lem:welfare-base}(b), and the
charging argument then bounds the approximation ratio by the
square root of \(\sum_j\nu_j^2\). Keeping this sum of order \(m\)
requires that a winner be not much larger than the bidders it
displaces, which is what thinning at the matching scale
enforces.

\subsection{Implementation and symmetry}
\label{sec:rand-implementation}

\begin{proposition}\label{prop:implementation}
The mechanism \(\mathcal A_{\mathcal K}\) is anonymous in
utilities and symmetric across goods. For \(n\) submitted
reports, it uses \(O(n^2\log n+n^2m)\) arithmetic operations.
\end{proposition}

\begin{proof}
The priority order is uniformly random on the participating
identities and independent of their reports, so relabelling
bidders does not change the distribution of outcomes. All other
operations depend on goods only through cardinalities,
intersections, and independent identical survival
probabilities. The mechanism is therefore anonymous in utilities
and symmetric across goods.

Only one scale is executed. One greedy run takes
\(O(n\log n+nm)\) elementary operations with explicit set
representations, and computing all needed critical weights by
deletion runs takes \(O(n^2\log n+n^2m)\) operations.
\end{proof}

\paragraph{Exact arithmetic.}
The payment~\eqref{eq:actual-payment} is a continuous random
variable whenever \(\xi C>\kappa_i\). If exact finite-precision
output is required, a served bidder of \(\mathcal A_k\) can
instead be charged its conditional expected payment
\(\bar P_i/(q_k(S_i)r_i)\), where, for \(C>0\),
\begin{equation}\label{eq:retention-payment}
r_i=\min\left\{1,\frac{w_i}{C}\right\},
\qquad
\bar P_i
=\E_\xi\bigl[\mathbf{1}\{\xi C\le w_i\}\max\{\kappa_i,\xi C\}\bigr]
=\frac{\min\{w_i,C\}^2+\kappa_i^2}{2C},
\end{equation}
and the payment is zero if \(C=0\). This is the payment that the
characterization of \citet{myerson1981optimal} associates with
the service probability \(r_i\). The allocation and all expected
payments are unchanged, so the welfare guarantee, truthfulness
in expectation, and false-name-proofness in expectation are
preserved. Since \(\kappa_i\le\min\{w_i,C\}\) by
Lemma~\ref{lem:greedy}(b), we have \(\bar P_i\le w_ir_i\); this
payment is therefore at most \(b_i\), and the mechanism remains
ex-post individually rational. It is no longer universally
truthful, because a winner's payment then depends on its own
bid. For rational reports, \(q_k(S)\) has denominator
\((2k)^{|S|}\), so weights, service probabilities, and these
payments have polynomial bit length. The reserve lottery can be
implemented by sampling among the at most \(|G|+1\) intervals
determined by the rational thresholds \(r_i\), and standard
exact rational sampling has expected polynomial bit complexity.

\paragraph{Scope of the result.}
The false-name guarantee is in expectation, with the bidder
liable for the payments of all its aliases. The deletion step
does not impose an extension of independence of irrelevant
goods: adding a privately desired good changes the probability
that the entire target survives. Theorem~\ref{thm:thinning} uses
nothing about greedy or about packing beyond exactness, and
Theorem~\ref{thm:reserve} isolates the single inequality that a
reserve law must satisfy; we expect both to be useful beyond the
present mechanism.

\bibliographystyle{ACM-Reference-Format}
\bibliography{sample-bibliography}
\end{document}
